%\documentclass[letterpaper,12pt]{article}
\documentclass[preview]{standalone}

\usepackage{tabularx}
\usepackage{amsmath}
\usepackage{graphicx}
\usepackage[margin=1in,letterpaper]{geometry}
\usepackage{cite}
\usepackage[final]{hyperref}
\usepackage{datatool}
\usepackage{caption}
\usepackage{float}
\usepackage{gensymb}

\hypersetup{
        colorlinks=true,
        linkcolor=blue,
        citecolor=blue,
        filecolor=magenta,
        urlcolor=blue
}
\begin{document}

\setcounter{table}{2}
\begin{table}[h!]
\centering
\label{tab:fits}
\renewcommand{\arraystretch}{1.2}
\begin{tabular}{|c|c|c|c|c|c|}
\hline
R $\pm 1$ $\Omega$ & $A$ | V & $B$ | V & $\lambda$ & $\omega$ | rad/s & $\phi$ | $^\circ$ \\
\hline
839 &
$2.625 \pm 0.004$ &
$-2384 \pm 16$ &
$1595 \pm 10$ &
-- & -- \\

10   & 
$4.633 \pm 0.005$ &
-- &
$179.5 \pm 0.3$ &
$4455.8 \pm 0.3$ &
$0.917 \pm 0.001$ \\

500   & 
$2.949 \pm 0.005$ &
-- &
$2459 \pm 11$ &
$3541 \pm 15$ &
$0.401 \pm 0.004$ \\

1000   &
$2.645 \pm 0.004$ &
$-1930 \pm 14$ &
$1472 \pm 7$ &
-- & -- \\

3500   &
$-0.044 \pm 0.006$ &
$2.82 \pm 0.01$ &
wtf &
-- \\
\hline
\end{tabular}
\caption{Paramètres ajustés des différents régimes d'amortissement}
\end{table}
% %%
% \setcounter{table}{3}
% \begin{table}[h!]
% \centering
% \label{tab:fits}
% \renewcommand{\arraystretch}{1.2}
% \begin{tabular}{|c|c|c|c|c|c|}
% \hline
% Noyau & $A$ | V & $\lambda$ & $\omega$ | rad/s & T | s \\
% \hline
% Acier feuilleté     &
% $4.695 \pm 0.005$ &
% $111.1 \pm 0.2$ &
% $3376.4 \pm 0.2$ &
% $0.0018609 \pm 1 \cdot 10^{-7}$ \\

% Acier     &
% $4.35 \pm 0.01$ &
% $1179 \pm 5$ &
% $4512 \pm 5$ &
% $0.0013925 \pm 1 \cdot 10^{-6}$ \\

% Air     &
% $4.618 \pm 0.005$ &
% $308.8 \pm 0.5$ &
% $7200.6 \pm 0.5$ &
% $0.0008726 \pm 6 \cdot 10^{-8}$ \\

% Aluminium     &
% $4.501 \pm 0.007$ &
% $910 \pm 2$ &
% $8894 \pm 2$ &
% $0.0007064 \pm 2 \cdot 10^{-7}$ \\
% \hline
% \end{tabular}
% \caption{Paramètres ajustés des oscillations avec le pont RLC et R = $10 \Omega$}
% \end{table}

% % %
% \setcounter{table}{4}
% \begin{table}[h!]
% \centering
% \label{tab:fits}
% \renewcommand{\arraystretch}{1.2}
% \begin{tabular}{|c|c|c|c|c|}
% \hline
% Noyau & $A$ | V & $\lambda$ & $\omega$ | rad/s & T | s \\
% \hline
% Acier     &
% $4.597 \pm 0.005$ &
% $263.6 \pm 0.4$ &
% $5362.9 \pm 0.4$ &
% $0.0011716 \pm 9 \cdot 10^{-8}$ \\

% Acier feuilleté     &
% $4.565 \pm 0.005$ &
% $187.4 \pm 0.3$ &
% $3381.8 \pm 0.3$ &
% $0.0018579 \pm 2 \cdot 10^{-7}$ \\

% Air     &
% $3.75 \pm 0.01$ &
% $1276 \pm 6$ &
% $7014 \pm 6$ &
% $0.0008958 \pm 8 \cdot 10^{-7}$ \\

% Aluminium    &
% $4.212 \pm 0.006$ &
% $1173 \pm 3$ &
% $8875 \pm 3$ &
% $0.0007079 \pm 2 \cdot 10^{-7}$ \\
% \hline
% \end{tabular}
% \caption{Paramètres ajustés des oscillations avec la boite à décade et R = $10 \Omega$}
% \end{table}


%\[
%T = \frac{2\pi}{\omega}
%\]


%\[
%\Delta T = \left| \frac{dT}{d\omega} \right| \Delta\omega
%\]

%\[
%\frac{dT}{d\omega} = \frac{d}{d\omega} \left( \frac{2\pi}{\omega} \right) = -\frac{2\pi}{\omega^2}
%\]

%\[
%\left| \frac{dT}{d\omega} \right| = \frac{2\pi}{\omega^2}
%\]


%\[
%\Delta T = \frac{2\pi}{\omega^2} \Delta\omega
%\]



\end{document}

